% c10 原书 p.117–132（PDF p.124–139）
\section{上同调}
\label{sec:2}

\subsection{导出函子}
\label{sec:2.1}

设 $\cC$ 是阿贝尔范畴。$\cC$ 中的\emph{复形}\idx{fuxing}{复形}是 $\cC$ 中的一列态射
\[
\cdots\to A^i\xrightarrow{d^i}A^{i+1}\xrightarrow{d^{i+1}}A^{i+2}\to\cdots,
\]
使得对每个 $i$ 都有 $d^{i+1}d^i=0$。态射 $d^i$ 称为该复形的\emph{上边缘态射}\idx{shangbianyuan}{上边缘态射}。我们常把上述复形记为 $A^\cdot$。若 $A^i$ 只在 $i$ 的某个范围内给定，则对其余的 $i$ 我们约定 $A^i=0$。

设 $A^\cdot$ 与 $B^\cdot$ 是两个复形。复形的态射 $f^\cdot\colon A^\cdot\to B^\cdot$ 由态射 $f^i\colon A^i\to B^i$（$i\in\bZ$）组成，使得下图对每个 $i$ 都交换：
\[\begin{tikzcd}[column sep=3.2em,row sep=1.8em]
A^i \arrow[r,"d^i"] \arrow[d,"f^i"'] & A^{i+1} \arrow[d,"f^{i+1}"] \\
B^i \arrow[r,"d^i"'] & B^{i+1}
\end{tikzcd}\]
我们按显然的方式定义态射的复合。于是得到 $\cC$ 中复形的范畴，它是阿贝尔范畴。

分别定义复形 $A^\cdot$ 的第 $i$ 个\emph{上闭链对象}\idx{shangbilian}{上闭链对象} $Z^i(A^\cdot)$、第 $i$ 个上边缘对象 $B^i(A^\cdot)$、以及第 $i$ 个\emph{上同调对象}\idx{shangtongdiaoduixiang}{上同调对象} $H^i(A^\cdot)$ 为
\[
\begin{aligned}
Z^i(A^\cdot)&=\Ker d^i,\\
B^i(A^\cdot)&=\Image d^{i-1},\\
H^i(A^\cdot)&=Z^i(A^\cdot)/B^i(A^\cdot).
\end{aligned}
\]
\begin{tnote}
上面第二个公式中 $B^i(A^\cdot)=\Image d^{i-1}$ 用的是 $d^{i-1}$ 的像而不是 $d^i$ 的像，指标比上闭链的 $Z^i=\Ker d^i$ 低一位，这是初学者最容易弄错的地方。[依据：原书 p.117 第 3 段公式]
\end{tnote}
若 $H^i(A^\cdot)=0$ 对所有 $i$ 成立，则称复形 $A^\cdot$ 是\emph{无环的}\idx{wuhuan}{无环复形}。

复形的每个态射 $f^\cdot\colon A^\cdot\to B^\cdot$ 都在上同调对象上诱导态射 $H^i(A^\cdot)\to H^i(B^\cdot)$。若存在态射 $h^i\colon A^i\to B^{i-1}$（$i\in\bZ$）使得
\[
f^i-g^i=h^{i+1}d^i+d^{i-1}h^i,
\]
则称两个态射 $f^\cdot,g^\cdot\colon A^\cdot\to B^\cdot$ 是\emph{同伦}\idx{tonglun}{同伦}的。

于是 $f^\cdot$ 与 $g^\cdot$ 在上同调对象上诱导相同的态射。若 $f^\cdot$ 与 $g^\cdot$ 同伦，我们记 $f^\cdot\sim g^\cdot$。若存在态射 $f^\cdot\colon A^\cdot\to B^\cdot$ 与 $g^\cdot\colon B^\cdot\to A^\cdot$，使 $g^\cdot f^\cdot\sim\id_{A^\cdot}$ 且 $f^\cdot g^\cdot\sim\id_{B^\cdot}$，则称两个复形 $A^\cdot$ 与 $B^\cdot$ 是\emph{同伦等价}\idx{tonglundengjia}{同伦等价}的。此时 $f^\cdot$ 与 $g^\cdot$ 在上同调对象上诱导同构。

复形态射的序列
\[
A^\cdot\to B^\cdot\to C^\cdot
\]
称为\emph{正合}的，若
\[
A^i\to B^i\to C^i
\]
对每个 $i$ 都正合。给定复形的短正合序列
\[
0\to A^\cdot\to B^\cdot\to C^\cdot\to 0,
\]
对每个 $i$ 可定义典范态射 $\delta^i\colon H^i(C^\cdot)\to H^{i+1}(A^\cdot)$，使得下述序列正合：
\[
\cdots\xrightarrow{\delta^{i-1}}H^i(A^\cdot)\to H^i(B^\cdot)\to H^i(C^\cdot)\xrightarrow{\delta^i}H^{i+1}(A^\cdot)\to\cdots.
\]
这个序列称为与上述复形短正合序列相伴的上同调对象的\emph{长正合序列}。
\begin{tnote}
连接同态 $\delta^i$ 的构造原文只写了结论：取 $c\in Z^i(C^\cdot)$，把它提升为 $b\in B^i$，$\dd b$ 落在 $A^{i+1}$ 中（因为它在 $C^{i+1}$ 中的像是 $0$），$\delta^i$ 就取 $\dd b$ 所在的类；良定义由 $A^i\to B^i$ 的单性与图追踪给出。[依据：原书 p.118 第 1 段；细节为标准同调代数，如 \S 2.1 中复形短正合序列的蛇引理]
\end{tnote}

设 $F\colon\cC\to\cD$（resp.\ $G\colon\cC\to\cD$）是阿贝尔范畴之间的加性共变（resp.\ 反变）函子。若对 $\cC$ 中任意短正合序列
\[
0\to A\to B\to C\to 0,
\]
序列
\[
0\to F(A)\to F(B)\to F(C)\quad(\text{resp. } 0\to G(C)\to G(B)\to G(A))
\]
正合，则称 $F$（resp.\ $G$）是\emph{左正合}的。若
\[
F(A)\to F(B)\to F(C)\to 0\quad(\text{resp. } G(C)\to G(B)\to G(A)\to 0)
\]
正合，则称 $F$（resp.\ $G$）是\emph{右正合}的。若 $F$（resp.\ $G$）既左正合又右正合，则称它是\emph{正合}的。这等价于说，$F$（resp.\ $G$）把 $\cC$ 中任意正合序列变为 $\cD$ 中的正合序列。

\noindent\textbf{例}
\begin{enumerate}[label=\arabic*.,leftmargin=2.4em,itemsep=0.2em]
  \item 对任意阿贝尔范畴 $\cC$ 与 $\cC$ 中任意对象 $A$，函子 $\Hom(A,-)$ 是从 $\cC$ 到阿贝尔群范畴的共变左正合函子，而函子 $\Hom(-,A)$ 是反变左正合函子。
  \item 从拓扑空间 $X$ 上阿贝尔群层范畴到阿贝尔群范畴的函子 $\cF\mapsto\cF(X)$ 是左正合的。
  \item 设 $f\colon X\to Y$ 是拓扑空间之间的连续映射。则 $\cF\mapsto f_*\cF$ 是从 $X$ 上阿贝尔群层范畴到 $Y$ 上阿贝尔群层范畴的左正合函子。
  \item 设 $(X,\cO_X)$ 是环层空间，$\cF$ 是 $\cO_X$-模。则 $\cO_X$-模范畴上的函子 $\cF\tensor_{\cO_X}-$ 是右正合的。
\end{enumerate}

阿贝尔范畴 $\cC$ 中的对象 $I$ 称为\emph{内射对象}\idx{nesheduixiang}{内射对象}，若 $\Hom(-,I)$ 是正合函子。这等价于说：对 $\cC$ 中任意对象 $A$ 与 $A$ 的任意子对象 $B$，每个态射 $B\to I$ 都能延拓为态射 $A\to I$。$\cC$ 中的对象 $P$ 称为\emph{投射对象}\idx{tousheduixiang}{投射对象}，若 $\Hom(P,-)$ 是正合函子。这等价于说：对任意对象 $A$ 与 $A$ 的任意商对象 $B$，每个态射 $P\to B$ 都能提升为态射 $P\to A$。注意，一个对象是内射对象，当且仅当它在对偶范畴中是投射对象。因此，若某个断言对内射对象成立，则其对偶断言对投射对象成立。在下文中，我们主要对内射对象陈述结论。
\begin{tnote}
「对子对象延拓」这一判据在一般阿贝尔范畴中即为 $\Hom(-,I)$ 正合的等价形式；在 $R$-模范畴中它退化为只需对理想 $I\subset R$ 检验的 Baer 判据。原书没有点出这层联系，但对验证具体对象是否内射时很有用。[来源：nLab "injective object", https://ncatlab.org/nlab/show/injective+object（访问日期 2026-11-14）]
\end{tnote}

设
\[
0\to A\xrightarrow{i}B\to C\to 0
\]
是短正合序列。若 $A$ 是内射对象，则该序列分裂。事实上，$\id_A$ 可延拓为态射 $\pi\colon B\to A$，我们有 $\pi i=\id_A$。我们的断言随即由 \ref{it:1.1.3} 推出。

对象 $A$ 的\emph{内射分解}\idx{neshefenjie}{内射分解}是形如
\[
0\to I^0\to I^1\to\cdots
\]
的复形，连同单态射 $A\to I^0$，使得每个 $I^i$ 都是内射对象且序列
\[
0\to A\to I^0\to I^1\to\cdots
\]
正合。$A$ 的\emph{投射分解}\idx{toushefenjie}{投射分解}是形如
\[
\cdots\to P^{-1}\to P^0\to 0
\]
的复形，连同满态射 $P^0\to A$，使得每个 $P^i$ 都是投射对象且序列
\[
\cdots\to P^{-1}\to P^0\to A\to 0
\]
正合。若对 $\cC$ 中每个对象 $A$，都存在从 $A$ 到 $\cC$ 中内射对象的单态射 $A\to I$，则称范畴 $\cC$ \emph{有足够内射对象}。若对 $\cC$ 中每个对象 $A$，都存在从投射对象 $P$ 到 $A$ 的满态射 $P\to A$，则称 $\cC$ \emph{有足够投射对象}。

\begin{proposition}\label{it:2.1.1}
设 $\cC$ 是阿贝尔范畴。
\begin{enumerate}[label=(\roman*),leftmargin=2.4em]
  \item 若 $\cC$ 有足够内射对象，则 $\cC$ 中每个对象都有内射分解。
  \item 设
  \[
  0\to A\to A^0\to A^1\to\cdots,
  \]
  \[
  0\to B\to I^0\to I^1\to\cdots
  \]
  是两个复形，使得第一个序列正合且每个 $I^i$ 都是内射对象。则对任意态射 $f\colon A\to B$，存在复形态射 $f^\cdot\colon A^\cdot\to I^\cdot$，使下图交换：
  \[\begin{tikzcd}[column sep=2.6em,row sep=1.8em]
  0 \arrow[r] & A \arrow[r] \arrow[d,"f"] & A^0 \arrow[r] \arrow[d,"f^0"] & A^1 \arrow[r] \arrow[d,"f^1"] & \cdots \\
  0 \arrow[r] & B \arrow[r] & I^0 \arrow[r] & I^1 \arrow[r] & \cdots
  \end{tikzcd}\]
  我们称这样的 $f^\cdot$ 为延拓 $f$ 的态射。若 $g^\cdot\colon A^\cdot\to I^\cdot$ 也是延拓 $f$ 的态射，则 $f^\cdot$ 与 $g^\cdot$ 同伦。
  \item 对象 $A$ 的任意两个内射分解都同伦等价。
  \item 设
  \[
  0\to A\to B\to C\to 0
  \]
  是 $\cC$ 中的正合序列，$I^\cdot$ 与 $J^\cdot$ 分别是 $A$ 与 $C$ 的内射分解。则存在正合序列
  \[
  0\to B\to I^0\oplus J^0\to I^1\oplus J^1\to\cdots
  \]
  使下图交换：
  \[\begin{tikzcd}[column sep=2.4em,row sep=1.7em]
  & 0 \arrow[d] & 0 \arrow[d] & 0 \arrow[d] & \\
  0 \arrow[r] & A \arrow[r] \arrow[d] & B \arrow[r] \arrow[d] & C \arrow[r] \arrow[d] & 0 \\
  0 \arrow[r] & I^0 \arrow[r] \arrow[d] & I^0\oplus J^0 \arrow[r] \arrow[d] & J^0 \arrow[r] \arrow[d] & 0 \\
  0 \arrow[r] & I^1 \arrow[r] \arrow[d] & I^1\oplus J^1 \arrow[r] \arrow[d] & J^1 \arrow[r] \arrow[d] & 0 \\
  & \vdots & \vdots & \vdots &
  \end{tikzcd}\]
  其中态射 $I^\cdot\to I^\cdot\oplus J^\cdot$ 与 $I^\cdot\oplus J^\cdot\to J^\cdot$ 是典范的。
\end{enumerate}
\end{proposition}

\begin{proof}
(i) 设 $A$ 是 $\cC$ 中对象。选取单态射 $A\to I^0$ 使 $I^0$ 是内射对象。于是我们有正合序列
\[
0\to A\to I^0.
\]
假设已找到正合序列
\[
0\to A\to I^0\to\cdots\to I^{n-1}.
\]
选取单态射
\[
\coker(I^{n-2}\to I^{n-1})\to I^n
\]
使 $I^n$ 是内射对象。（当 $n=1$ 时，我们取 $I^{n-2}=A$。）于是我们有正合序列
\[
0\to A\to I^0\to\cdots\to I^n.
\]
如此便得到 $A$ 的内射分解。

(ii) 因 $I^0$ 是内射对象且 $A$ 是 $A^0$ 的子对象，复合
\[
A\xrightarrow{f}B\to I^0
\]
可延拓为态射 $f^0\colon A^0\to I^0$。下图交换：
\[\begin{tikzcd}[column sep=2.6em,row sep=1.8em]
0 \arrow[r] & A \arrow[r] \arrow[d,"f"] & A^0 \arrow[d,"f^0"] \\
0 \arrow[r] & B \arrow[r] & I^0
\end{tikzcd}\]
假设对使 $i<n$ 的每个 $i$ 都已定义 $f^i$，且下图交换：
\[\begin{tikzcd}[column sep=2.2em,row sep=1.8em]
0 \arrow[r] & A \arrow[r] \arrow[d,"f"] & A^0 \arrow[r] \arrow[d,"f^0"] & A^1 \arrow[r] \arrow[d,"f^1"] & \cdots \arrow[r] & A^{n-1} \arrow[d,"f^{n-1}"] \\
0 \arrow[r] & B \arrow[r] & I^0 \arrow[r] & I^1 \arrow[r] & \cdots \arrow[r] & I^{n-1}
\end{tikzcd}\]
注意复合
\[
A^{n-1}\xrightarrow{f^{n-1}}I^{n-1}\to I^n
\]
在 $\Image(A^{n-2}\to A^{n-1})$ 上为零，因此它诱导态射
\[
\coker(A^{n-2}\to A^{n-1})\to I^n.
\]
因 $I^n$ 是内射对象且 $\coker(A^{n-2}\to A^{n-1})$ 是 $A^n$ 的子对象，此态射可延拓为态射 $f^n\colon A^n\to I^n$。下图交换：
\[\begin{tikzcd}[column sep=2.2em,row sep=1.8em]
0 \arrow[r] & A \arrow[r] \arrow[d,"f"] & A^0 \arrow[r] \arrow[d,"f^0"] & A^1 \arrow[r] \arrow[d,"f^1"] & \cdots \arrow[r] & A^n \arrow[d,"f^n"] \\
0 \arrow[r] & B \arrow[r] & I^0 \arrow[r] & I^1 \arrow[r] & \cdots \arrow[r] & I^n
\end{tikzcd}\]
如此便得到延拓 $f\colon A\to B$ 的态射 $f^\cdot\colon A^\cdot\to I^\cdot$。

设 $g^\cdot\colon A^\cdot\to I^\cdot$ 也是延拓 $f$ 的态射。令 $I^{-1}=0$，并定义 $h^0\colon A^0\to I^{-1}$ 为零态射。假设对每个 $i\le n$ 都已定义 $h^i\colon A^i\to I^{i-1}$，使得
\[
h^{i+1}d^i+d^{i-1}h^i=f^i-g^i
\]
对每个 $i\le n-1$ 成立。我们有
\[
\begin{aligned}
&(f^n-g^n-d^{n-1}h^n)d^{n-1}\\
&\quad=d^{n-1}(f^{n-1}-g^{n-1})-d^{n-1}h^n d^{n-1}\\
&\quad=d^{n-1}(h^n d^{n-1}+d^{n-2}h^{n-1})-d^{n-1}h^n d^{n-1}\\
&\quad=0.
\end{aligned}
\]
于是 $f^n-g^n-d^{n-1}h^n$ 在 $\Image(A^{n-1}\to A^n)$ 上为零，从而诱导从 $\coker(A^{n-1}\to A^n)$ 到 $I^n$ 的态射。但 $\coker(A^{n-1}\to A^n)$ 是 $A^{n+1}$ 的子对象，而 $I^n$ 是内射对象，故此态射可延拓为态射 $h^{n+1}\colon A^{n+1}\to I^n$。于是我们有
\[
h^{n+1}d^n+d^{n-1}h^n=f^n-g^n.
\]

(iii) 设 $I^\cdot$ 与 $J^\cdot$ 是 $A$ 的两个内射分解。由 (ii)，存在延拓 $\id_A$ 的态射 $f^\cdot\colon I^\cdot\to J^\cdot$，也存在延拓 $\id_A$ 的态射 $g^\cdot\colon J^\cdot\to I^\cdot$。注意 $g^\cdot f^\cdot\colon I^\cdot\to I^\cdot$ 与 $\id_{I^\cdot}$ 都是延拓 $\id_A$ 的态射。于是由 (ii)，$g^\cdot f^\cdot\sim\id_{I^\cdot}$。同理 $f^\cdot g^\cdot\sim\id_{J^\cdot}$。所以 $I^\cdot$ 与 $J^\cdot$ 同伦等价。

(iv) 因 $I^0$ 是内射对象且 $A$ 是 $B$ 的子对象，态射 $A\to I^0$ 可延拓为态射 $B\to I^0$。设
\[
B\to I^0\oplus J^0
\]
是这样的态射：它与投影 $I^0\oplus J^0\to I^0$ 及 $I^0\oplus J^0\to J^0$ 的复合分别是 $B\to I^0$ 与 $B\to C\to J^0$。我们有交换图
\[\begin{tikzcd}[column sep=2.4em,row sep=1.8em]
0 \arrow[r] & A \arrow[r] \arrow[d] & B \arrow[r] \arrow[d] & C \arrow[r] \arrow[d] & 0 \\
0 \arrow[r] & I^0 \arrow[r] & I^0\oplus J^0 \arrow[r] & J^0 \arrow[r] & 0.
\end{tikzcd}\]
第一个与第三个竖直箭头是单态射。故 $B\to I^0\oplus J^0$ 是单态射。假设对使 $i\le n$ 的每个 $i$ 都已定义态射 $I^{i-1}\oplus J^{i-1}\to I^i\oplus J^i$，使得下图交换且所有列正合：
\[\begin{tikzcd}[column sep=2.4em,row sep=1.5em]
& 0 \arrow[d] & 0 \arrow[d] & 0 \arrow[d] & \\
0 \arrow[r] & A \arrow[r] \arrow[d] & B \arrow[r] \arrow[d] & C \arrow[r] \arrow[d] & 0 \\
0 \arrow[r] & I^0 \arrow[r] \arrow[d] & I^0\oplus J^0 \arrow[r] \arrow[d] & J^0 \arrow[r] \arrow[d] & 0 \\
& \vdots \arrow[d] & \vdots \arrow[d] & \vdots \arrow[d] & \\
0 \arrow[r] & I^n \arrow[r] & I^n\oplus J^n \arrow[r] & J^n \arrow[r] & 0
\end{tikzcd}\]
则序列
\[
0\to\coker(I^{n-1}\to I^n)\to\coker(I^{n-1}\oplus J^{n-1}\to I^n\oplus J^n)\to\coker(J^{n-1}\to J^n)\to 0
\]
正合。事实上，这是与复形的短正合序列
\[
0\to I^\cdot\to I^\cdot\oplus J^\cdot\to J^\cdot\to 0
\]
相伴的上同调对象的长正合序列的一部分，其中 $I^\cdot$ 是复形
\[
0\to A\to I^0\to\cdots\to I^n\to 0,
\]
$J^\cdot$ 与 $I^\cdot\oplus J^\cdot$ 类似。因 $I^{n+1}$ 是内射对象，我们可以把态射 $\coker(I^{n-1}\to I^n)\to I^{n+1}$ 延拓为态射
\[
\coker(I^{n-1}\oplus J^{n-1}\to I^n\oplus J^n)\to I^{n+1}.
\]
此态射连同复合
\[
\coker(I^{n-1}\oplus J^{n-1}\to I^n\oplus J^n)\to\coker(J^{n-1}\to J^n)\to J^{n+1}
\]
定义了态射
\[
\coker(I^{n-1}\oplus J^{n-1}\to I^n\oplus J^n)\to I^{n+1}\oplus J^{n+1},
\]
它必是单态射。于是下图交换且所有列正合：
\[\begin{tikzcd}[column sep=2.4em,row sep=1.5em]
& 0 \arrow[d] & 0 \arrow[d] & 0 \arrow[d] & \\
0 \arrow[r] & A \arrow[r] \arrow[d] & B \arrow[r] \arrow[d] & C \arrow[r] \arrow[d] & 0 \\
0 \arrow[r] & I^0 \arrow[r] \arrow[d] & I^0\oplus J^0 \arrow[r] \arrow[d] & J^0 \arrow[r] \arrow[d] & 0 \\
& \vdots \arrow[d] & \vdots \arrow[d] & \vdots \arrow[d] & \\
0 \arrow[r] & I^{n+1} \arrow[r] & I^{n+1}\oplus J^{n+1} \arrow[r] & J^{n+1} \arrow[r] & 0
\end{tikzcd}\]
\end{proof}

设 $F\colon\cC\to\cD$ 是阿贝尔范畴之间的左正合共变函子。设 $\cC$ 有足够内射对象。我们如下定义 $F$ 的\emph{右导出函子}\idx{youdaochuhanzi}{右导出函子} $\R^iF$（$i=0,1,\dots$）：对 $\cC$ 中每个对象 $A$，一次性选定 $A$ 的一个内射分解 $I_A^\cdot$。对每个 $i$，定义 $\R^iF(A)$ 为复形 $F(I_A^\cdot)$ 的第 $i$ 个上同调对象。给定 $\cC$ 中的态射 $A\to B$，存在延拓 $A\to B$ 的复形态射 $I_A^\cdot\to I_B^\cdot$，由 \ref{it:2.1.1}(ii) 它在同伦意义下唯一。它诱导良定义的态射
\[
\R^iF(A)=H^i(F(I_A^\cdot))\to H^i(F(I_B^\cdot))=\R^iF(B).
\]
由 \ref{it:2.1.1}(iii)，函子 $\R^iF$ 的同构类不依赖于内射分解的选取。利用 $F$ 左正合的假设，可以验证 $\R^0F=F$。对 $\cC$ 中每个内射对象 $I$，我们有 $\R^iF(I)=0$ 对所有 $i\ge1$ 成立。

给定短正合序列
\[
0\to A\to B\to C\to 0,
\]
由 \ref{it:2.1.1}(iv)，可找到复形的短正合序列
\[
0\to I_A^\cdot\to I_B^\cdot\to I_C^\cdot\to 0,
\]
使得 $I_A^\cdot$、$I_B^\cdot$、$I_C^\cdot$ 分别是 $A$、$B$、$C$ 的内射分解，且对每个 $i$ 有
\[
I_B^i\cong I_A^i\oplus I_C^i.
\]
于是
\[
0\to F(I_A^\cdot)\to F(I_B^\cdot)\to F(I_C^\cdot)\to 0
\]
是复形的（分裂的）短正合序列。因此对每个 $i$ 存在态射
\[
\delta^i\colon\R^iF(C)\to\R^{i+1}F(A),
\]
使得序列
\[
\cdots\xrightarrow{\delta^{i-1}}\R^iF(A)\to\R^iF(B)\to\R^iF(C)\xrightarrow{\delta^i}\R^{i+1}F(A)\to\cdots
\]
正合。我们称此序列为与短正合序列
\[
0\to A\to B\to C\to 0
\]
相伴的 $\R^iF$ 的\emph{长正合序列}。
\begin{tnote}
$\R^iF$ 与内射分解的选取无关，关键在于 2.1.1(ii)(iii)：任意两个内射分解同伦等价，而同伦的复形态射在 $F$ 之后仍同伦（$F$ 是加性函子），故诱导相同的上同调态射。$\R^0F=F$ 则用到 $F$ 左正合，这一点在取出 $H^0$ 时不能省。[依据：原书 p.124 第 1 段与 p.121–122 的证明]
\end{tnote}

给定短正合序列的态射
\[\begin{tikzcd}[column sep=2.2em,row sep=1.5em]
0 \arrow[r] & A \arrow[r] \arrow[d] & B \arrow[r] \arrow[d] & C \arrow[r] \arrow[d] & 0 \\
0 \arrow[r] & A' \arrow[r] & B' \arrow[r] & C' \arrow[r] & 0,
\end{tikzcd}\]
则下图对每个 $i$ 交换：
\[\begin{tikzcd}[column sep=3.2em,row sep=1.8em]
\R^iF(C) \arrow[r,"\delta^i"] \arrow[d] & \R^{i+1}F(A) \arrow[d] \\
\R^iF(C') \arrow[r,"\delta^i"'] & \R^{i+1}F(A').
\end{tikzcd}\]

$\cC$ 中的对象 $J$ 称为 \emph{$F$-无环}的，若 $\R^iF(J)=0$ 对所有 $i\ge1$ 成立。对象 $A$ 的 \emph{$F$-无环分解}是形如
\[
0\to J^0\to J^1\to\cdots
\]
的复形，连同单态射 $A\to J^0$，使得每个 $J^i$ 都是 $F$-无环的且序列
\[
0\to A\to J^0\to J^1\to\cdots
\]
正合。

\begin{proposition}\label{it:2.1.2}
设 $F\colon\cC\to\cD$ 是阿贝尔范畴之间的左正合共变函子，且 $\cC$ 有足够内射对象。设 $J^\cdot$ 是 $A$ 的 $F$-无环分解。则对每个 $i$ 有
\[
\R^iF(A)\cong H^i(F(J^\cdot)).
\]
\end{proposition}
\begin{proof}
由与短正合序列
\[
0\to A\to J^0\to Z^1(J^\cdot)\to 0
\]
相伴的 $\R^iF$ 的长正合序列，以及 $\R^iF(J^0)=0$（$i\ge1$），得到
\[
\R^{i-1}F(Z^1(J^\cdot))\cong\R^iF(A)
\]
对每个 $i\ge2$ 成立。由与短正合序列
\[
0\to Z^1(J^\cdot)\to J^1\to Z^2(J^\cdot)\to 0
\]
相伴的 $\R^iF$ 的长正合序列，以及 $\R^iF(J^1)=0$（$i\ge1$），得到
\[
\R^{i-2}F(Z^2(J^\cdot))\cong\R^{i-1}F(Z^1(J^\cdot))\cong\R^iF(A)
\]
对每个 $i\ge3$ 成立。如此继续，得到
\[
\R^iF(A)\cong\R^{i-1}F(Z^1(J^\cdot))\cong\cdots\cong\R^1F(Z^{i-1}(J^\cdot))
\]
对每个 $i\ge1$ 成立。由与短正合序列
\[
0\to Z^{i-1}(J^\cdot)\to J^{i-1}\to Z^i(J^\cdot)\to 0
\]
相伴的 $\R^iF$ 的长正合序列，以及 $\R^1F(J^{i-1})=0$，得到正合序列
\[
0\to F(Z^{i-1}(J^\cdot))\to F(J^{i-1})\to F(Z^i(J^\cdot))\to\R^1F(Z^{i-1}(J^\cdot))\to 0,
\]
从而
\[
\R^1F(Z^{i-1}(J^\cdot))\cong F(Z^i(J^\cdot))/\Image(F(J^{i-1})\to F(Z^i(J^\cdot))).
\]
我们有正合序列
\[
0\to Z^i(J^\cdot)\to J^i\to J^{i+1}.
\]
因 $F$ 左正合，我们有正合序列
\[
0\to F(Z^i(J^\cdot))\to F(J^i)\to F(J^{i+1}),
\]
从而
\[
F(Z^i(J^\cdot))\cong\Ker(F(J^i)\to F(J^{i+1})),
\]
\[
\Image(F(J^{i-1})\to F(Z^i(J^\cdot)))\cong\Image(F(J^{i-1})\to F(J^i)).
\]
于是
\[
\begin{aligned}
\R^1F(Z^{i-1}(J^\cdot))&\cong\Ker(F(J^i)\to F(J^{i+1}))/\Image(F(J^{i-1})\to F(J^i))\\
&\cong H^i(F(J^\cdot)).
\end{aligned}
\]
因此 $\R^iF(A)\cong H^i(F(J^\cdot))$ 对每个 $i\ge1$ 成立。利用 $F$ 左正合，可以证明 $\R^0F(A)\cong H^0(F(J^\cdot))$。
\end{proof}

\begin{proposition}\label{it:2.1.3}
设 $F\colon\cC\to\cD$ 是阿贝尔范畴之间的左正合共变函子，且 $\cC$ 有足够内射对象。设 $\cS$ 是一族对象，满足以下条件：
\begin{enumerate}[label=(\alph*),leftmargin=2.4em]
  \item 对 $\cC$ 中每个对象 $A$，存在从 $A$ 到 $\cS$ 中某个对象的单态射。
  \item 若 $\cC$ 中对象 $A$ 同构于 $\cS$ 中某个对象的直和项，则 $A$ 是 $\cS$ 中的对象。
  \item 设
  \[
  0\to A\to B\to C\to 0
  \]
  是 $\cC$ 中的正合序列。若 $A$ 与 $B$ 都是 $\cS$ 中的对象，则 $C$ 也是 $\cS$ 中的对象，且序列
  \[
  0\to F(A)\to F(B)\to F(C)\to 0
  \]
  正合。
\end{enumerate}
则 $\cC$ 中每个内射对象都是 $\cS$ 中的对象，且 $\cS$ 中每个对象都是 $F$-无环的。
\end{proposition}
\begin{proof}
设 $I$ 是 $\cC$ 中的内射对象，$I\to S$ 是从 $I$ 到 $\cS$ 中某对象 $S$ 的单态射。把此单态射的余核记为 $S/I$。于是我们有分裂正合序列
\[
0\to I\to S\to S/I\to 0.
\]
所以 $I$ 同构于 $S$ 的直和项。由 (b)，$I$ 是 $\cS$ 中的对象。

设 $S$ 是 $\cS$ 中的对象，$I^\cdot$ 是 $S$ 的内射分解。于是我们有正合序列
\[
\begin{aligned}
0&\to S\to I^0\to Z^1(I^\cdot)\to 0,\\
0&\to Z^1(I^\cdot)\to I^1\to Z^2(I^\cdot)\to 0,\\
0&\to Z^2(I^\cdot)\to I^2\to Z^3(I^\cdot)\to 0,\\
&\ \vdots
\end{aligned}
\]
由 (c)，这些短正合序列依次表明 $Z^1(I^\cdot),Z^2(I^\cdot),\dots$ 都是 $\cS$ 中的对象，且下述序列正合：
\[
\begin{aligned}
0&\to F(S)\to F(I^0)\to F(Z^1(I^\cdot))\to 0,\\
0&\to F(Z^1(I^\cdot))\to F(I^1)\to F(Z^2(I^\cdot))\to 0,\\
0&\to F(Z^2(I^\cdot))\to F(I^2)\to F(Z^3(I^\cdot))\to 0,\\
&\ \vdots
\end{aligned}
\]
这表明
\[
0\to F(S)\to F(I^0)\to F(I^1)\to F(I^2)\to\cdots
\]
正合。所以 $\R^iF(S)\cong H^i(F(I^\cdot))=0$ 对每个 $i\ge1$ 成立。因此 $\cS$ 中的对象都是 $F$-无环的。
\end{proof}

设 $\cC$ 与 $\cD$ 是阿贝尔范畴。从 $\cC$ 到 $\cD$ 的共变\emph{上同调函子}\idx{shangtongdiaohanzi}{上同调函子}是一族共变函子 $T=(T^i)$（$i=0,1,\dots$），连同对每个短正合序列
\[
0\to A\to B\to C\to 0
\]
给出的态射 $\delta^i\colon T^i(C)\to T^{i+1}(A)$，满足以下条件：
\begin{enumerate}[label=(\roman*),leftmargin=2.4em]
  \item 与上述短正合序列相伴，我们有长正合序列
  \[
  0\to T^0(A)\to T^0(B)\to T^0(C)\xrightarrow{\delta^0}T^1(A)\to\cdots.
  \]
  \item 对短正合序列的每个态射
  \[\begin{tikzcd}[column sep=2.2em,row sep=1.5em]
  0 \arrow[r] & A \arrow[r] \arrow[d] & B \arrow[r] \arrow[d] & C \arrow[r] \arrow[d] & 0 \\
  0 \arrow[r] & A' \arrow[r] & B' \arrow[r] & C' \arrow[r] & 0,
  \end{tikzcd}\]
  下图对每个 $i$ 交换：
  \[\begin{tikzcd}[column sep=3.2em,row sep=1.8em]
  T^i(C) \arrow[r,"\delta^i"] \arrow[d] & T^{i+1}(A) \arrow[d] \\
  T^i(C') \arrow[r,"\delta^i"'] & T^{i+1}(A').
  \end{tikzcd}\]
\end{enumerate}

若对每个上同调函子 $T'=(T'^i)$ 与每个自然变换 $f^0\colon T^0\to T'^0$，都存在唯一一族自然变换 $f^i\colon T^i\to T'^i$（$i\ge1$），使得对每个短正合序列
\[
0\to A\to B\to C\to 0,
\]
下图对每个 $i$ 交换，则称 $T=(T^i)$ 是\emph{泛}\idx{fanshangtongdiaohanzi}{泛上同调函子}的：
\[\begin{tikzcd}[column sep=3.6em,row sep=1.8em]
T^i(C) \arrow[r,"\delta^i"] \arrow[d,"f^i_C"'] & T^{i+1}(A) \arrow[d,"f^{i+1}_A"] \\
T'^i(C) \arrow[r,"\delta'^i"'] & T'^{i+1}(A).
\end{tikzcd}\]

函子 $F\colon\cC\to\cD$ 称为\emph{可抹的}\idx{kemohanzi}{可抹函子}，若对 $\cC$ 中任意对象 $A$，存在单态射 $u\colon A\to M$ 使 $F(u)=0$。若对 $\cC$ 中任意对象 $A$，存在满态射 $u\colon M\to A$ 使 $F(u)=0$，则称 $F$ 是\emph{可余抹的}\idx{keyumo}{可余抹函子}。
\begin{tnote}
可抹与可余抹是互相对偶的概念：把 $\cC$ 换成对偶范畴，单态射与满态射互换，箭头反向，可抹函子即变为可余抹函子。原书只对共变上同调函子陈述 2.1.4，反变情形可由对偶性得到。[依据：原书 p.128 定义段]
\end{tnote}

\begin{proposition}\label{it:2.1.4}
设 $T=(T^i)\colon\cC\to\cD$ 是阿贝尔范畴之间的共变上同调函子。若对每个 $i\ge1$，$T^i$ 都是可抹的，则 $T$ 是泛的。
\end{proposition}
\begin{proof}
设 $T'=(T'^i)$ 是上同调函子，$f^0\colon T^0\to T'^0$ 是自然变换。假设已证明：存在唯一一族自然变换 $f^i\colon T^i\to T'^i$（$0\le i\le n-1$），使得对任意短正合序列
\[
0\to A\to B\to C\to 0,
\]
下图对 $0\le i\le n-2$ 都交换：
\[\begin{tikzcd}[column sep=3.6em,row sep=1.8em]
T^i(C) \arrow[r,"\delta^i"] \arrow[d,"f^i_C"'] & T^{i+1}(A) \arrow[d,"f^{i+1}_A"] \\
T'^i(C) \arrow[r,"\delta'^i"'] & T'^{i+1}(A).
\end{tikzcd}\]
对 $\cC$ 中每个对象 $A$，选取单态射 $u\colon A\to M$ 使 $T^n(u)=0$，并令 $M/A$ 为其余核。我们有短正合序列
\[
0\to A\to M\to M/A\to 0.
\]
考虑下图：
\[\begin{tikzcd}[column sep=3.2em,row sep=2.2em]
T^{n-1}(M) \arrow[r] \arrow[d,"f^{n-1}_M"] & T^{n-1}(M/A) \arrow[r] \arrow[d,"f^{n-1}_{M/A}"] & T^n(A) \arrow[r] \arrow[d,dashed,"f^n_A"] & 0 \\
T'^{n-1}(M) \arrow[r] & T'^{n-1}(M/A) \arrow[r] & T'^n(A) &
\end{tikzcd}\]
注意两行都正合。因此存在唯一态射
\[
f^n_A\colon T^n(A)\to T'^n(A)
\]
使上图交换。

我们断言，态射 $f^n_A$ 不依赖于满足 $T^n(u)=0$ 的单态射 $u\colon A\to M$ 的选取。设 $v\colon A\to N$ 是另一个满足 $T^n(v)=0$ 的单态射，$L$ 为由 $-u\colon A\to M$ 与 $v\colon A\to N$ 定义的态射 $A\to M\oplus N$ 的余核。于是我们有交换图
\[\begin{tikzcd}[column sep=2.8em,row sep=1.8em]
A \arrow[r,"u"] \arrow[d,"v"'] \arrow[dr,"w"] & M \arrow[d] \\
N \arrow[r] & L,
\end{tikzcd}\]
且态射 $w\colon A\to L$ 是满足 $T^n(w)=0$ 的单态射。态射 $M\to L$ 诱导态射 $M/A\to L/A$。交换图
\[\begin{tikzcd}[column sep=3.6em,row sep=1.8em]
T^{n-1}(M/A) \arrow[r] \arrow[d,"f^{n-1}_{M/A}"'] & T^{n-1}(L/A) \arrow[d,"f^{n-1}_{L/A}"] \\
T'^{n-1}(M/A) \arrow[r] & T'^{n-1}(L/A)
\end{tikzcd}\]
说明：由单态射 $u$ 定义的态射 $f^n_A\colon T^n(A)\to T'^n(A)$ 与由 $w$ 定义的那个重合。同理，由单态射 $v$ 定义的 $f^n_A$ 也与由 $w$ 定义的那个重合。

设 $\alpha\colon A\to B$ 是 $\cC$ 中的态射。选取单态射 $u\colon A\to M$ 与 $v_0\colon B\to N_0$，使 $T^n(u)=0$ 且 $T^n(v_0)=0$。令 $N$ 为由 $-u\colon A\to M$ 与 $v_0\alpha\colon A\to N_0$ 定义的态射 $A\to M\oplus N_0$ 的余核。于是我们有交换图
\[\begin{tikzcd}[column sep=2.8em,row sep=1.8em]
A \arrow[r,"u"] \arrow[d,"\alpha"'] & M \arrow[d] \\
B \arrow[r,"v"'] & N,
\end{tikzcd}\]
其中 $u$ 与 $v$ 都是单态射且满足 $T^n(u)=0$ 与 $T^n(v)=0$。我们有交换图
\[\begin{tikzcd}[column sep=3.6em,row sep=1.8em]
T^{n-1}(M/A) \arrow[r] \arrow[d,"f^{n-1}_{M/A}"'] & T^{n-1}(N/B) \arrow[d,"f^{n-1}_{N/B}"] \\
T'^{n-1}(M/A) \arrow[r] & T'^{n-1}(N/B).
\end{tikzcd}\]
这说明下图交换：
\[\begin{tikzcd}[column sep=3.2em,row sep=1.8em]
T^n(A) \arrow[r] \arrow[d,"f^n_A"'] & T^n(B) \arrow[d,"f^n_B"] \\
T'^n(A) \arrow[r] & T'^n(B).
\end{tikzcd}\]
所以 $f^n$ 是自然变换。

给定短正合序列
\[
0\to A\to B\to C\to 0,
\]
选取单态射 $v\colon B\to M$ 使 $T^n(v)=0$。令 $u$ 为复合 $A\to B\xrightarrow{v}M$。则 $T^n(u)=0$。我们有交换图
\[\begin{tikzcd}[column sep=2.6em,row sep=1.6em]
0 \arrow[r] & A \arrow[r] \arrow[d,equal] & B \arrow[r] \arrow[d,"v"'] & C \arrow[r] \arrow[d] & 0 \\
0 \arrow[r] & A \arrow[r,"u"'] & M \arrow[r] & M/A \arrow[r] & 0.
\end{tikzcd}\]
它给出下述交换图：
\[\begin{tikzcd}[column sep=2.4em,row sep=3.0em]
T^{n-1}(C) \arrow[r] \arrow[d,"\delta^{n-1}"'] \arrow[rrr,bend left=26,"f^{n-1}_C"]
  & T^{n-1}(M/A) \arrow[r,"f^{n-1}_{M/A}"] \arrow[d,"\delta^{n-1}"']
  & T'^{n-1}(M/A) \arrow[d,"\delta^{n-1}"']
  & T'^{n-1}(C) \arrow[l] \arrow[d,"\delta^{n-1}"'] \\
T^n(A) \arrow[r,"\id_{T^n(A)}"]
  & T^n(A) \arrow[r,"f^n_A"]
  & T'^n(A)
  & T'^n(A) \arrow[l,"\id_{T'^n(A)}"']
\end{tikzcd}\]
这说明下图交换：
\[\begin{tikzcd}[column sep=4.2em,row sep=1.8em]
T^{n-1}(C) \arrow[r,"\delta^{n-1}"] \arrow[d,"f^{n-1}_C"'] & T^n(A) \arrow[d,"f^n_A"] \\
T'^{n-1}(C) \arrow[r,"\delta^{n-1}"'] & T'^n(A).
\end{tikzcd}\]
\end{proof}

\begin{proposition}\label{it:2.1.5}
设 $F\colon\cC\to\cD$ 是阿贝尔范畴之间的左正合共变函子，且 $\cC$ 有足够内射对象。则 $(\R^iF)$ 是泛上同调函子。反之，若 $T=(T^i)$ 是共变泛上同调函子，则 $T^0$ 是左正合的，且 $T^i$ 对每个 $i\ge0$ 都同构于 $\R^iT^0$。
\end{proposition}
\begin{proof}
由定义，$(\R^iF)$ 是上同调函子。对 $\cC$ 中每个对象 $A$，选取单态射 $u\colon A\to I$ 使 $I$ 是内射对象。因 $\R^iF(I)=0$ 对每个 $i\ge1$ 成立，我们有 $\R^iF(u)=0$ 对每个 $i\ge1$ 成立。这说明 $\R^iF$ 是可抹的。所以由 \ref{it:2.1.4}，$(\R^iF)$ 是泛的。

反之，设 $T=(T^i)$ 是共变泛上同调函子。则由上同调函子的定义，$T^0$ 左正合。$(T^i)$ 与 $(\R^iT^0)$ 都是泛上同调函子，且它们的第 $0$ 个分量相同。所以它们同构。
\end{proof}

我们看到，当 $\cC$ 有足够内射对象时，可用内射分解定义左正合共变函子 $F$ 的右导出函子 $\R^iF$，它们构成泛上同调函子。类似地，若 $\cC$ 有足够投射对象，我们可用投射分解定义右正合共变函子 $F$ 的\emph{左导出函子}\idx{zuodaochuhanzi}{左导出函子} $L_iF$。对任意反变右正合（resp.\ 左正合）函子，我们也可用内射（resp.\ 投射）分解定义它的左导出函子（resp.\ 右导出函子）。它们的泛性质留给读者陈述。

\begin{proposition}\label{it:2.1.6}
设 $(X,\cO_X)$ 是环层空间。则 $\cO_X$-模层构成的阿贝尔范畴有足够内射对象。
\end{proposition}
\begin{proof}
对 $X$ 的每个开子集 $U$，定义预层 $\cP_U$ 为
\[
\cP_U(V)=
\begin{cases}
\cO_X(V), & \text{若 } V\subset U,\\
\{0\}, & \text{若 } V\not\subset U.
\end{cases}
\]
设 $\cG_U$ 是预层 $\cP_U$ 的相伴层。给定 $\cO_X$-模 $\cF$ 与截面 $s\in\cF(U)$，定义态射
\[
\cP_U\to\cF
\]
如下：若 $V\subset U$，定义 $\cP_U(V)=\cO_X(V)\to\cF(V)$ 为 $a\mapsto as|_V$；若 $V\not\subset U$，定义 $\cP_U(V)=\{0\}\to\cF(V)$ 为零态射。态射 $\cP_U\to\cF$ 诱导态射 $\cG_U\to\cF$，我们称之为由截面 $s\in\cF(U)$ 定义的态射。令 $\cG=\bigoplus_U\cG_U$，其中直和取遍 $X$ 的所有开子集。则 $\cG$ 是 $\cO_X$-模。

我们先证明：若对 $\cG$ 的每个 $\cO_X$-子模 $\cG'$，每个从 $\cG'$ 到 $\cI$ 的态射都能延拓为从 $\cG$ 到 $\cI$ 的态射，则 $\cO_X$-模 $\cI$ 是内射的。给定 $\cO_X$-模 $\cF$、子模 $\cF'\subset\cF$ 与态射 $\phi'\colon\cF'\to\cI$，我们需要证明存在态射 $\phi\colon\cF\to\cI$ 使 $\phi|_{\cF'}=\phi'$。令 $\cS$ 为所有对 $(\cF'',\phi'')$ 构成的族，其中 $\cF''$ 是 $\cF$ 的含 $\cF'$ 的 $\cO_X$-子模，$\phi''\colon\cF''\to\cI$ 是满足 $\phi''|_{\cF'}=\phi'$ 的态射。由 Zorn 引理，$\cS$ 有极大元。为证明上述断言，只需证明 $\cS$ 中每个极大元 $(\cF'',\phi'')$ 都满足 $\cF''=\cF$。为简化记号，设 $(\cF',\phi')$ 是极大元，我们来证 $\cF'=\cF$。若 $\cF'\ne\cF$，则存在开子集 $W$ 与截面 $s\in\cF(W)-\cF'(W)$。定义态射
\[
\psi\colon\cG=\bigoplus_U\cG_U\to\cF,
\]
使得当 $U\ne W$ 时 $\psi|_{\cG_U}=0$，而 $\psi|_{\cG_W}\colon\cG_W\to\cF$ 是由截面 $s$ 定义的态射。令 $\cG'$ 为 $\cF'$ 在 $\psi$ 下的原像，$\psi'\colon\cG'\to\cF'$ 为 $\psi$ 诱导的态射。由假设，存在态射 $\theta\colon\cG\to\cI$ 使 $\theta|_{\cG'}=\phi'\psi'$。我们有正合序列
\[
0\to\cG'\xrightarrow{\alpha}\cF'\oplus\cG\xrightarrow{\beta}\cF,
\]
其中 $\alpha$ 与 $\beta$ 由
\[
\alpha(g')=(\psi'(g'),-g'),\qquad \beta(f',g)=f'+\psi(g)
\]
定义。

考虑态射
\[
\delta'\colon\cF'\oplus\cG\to\cI,\qquad (f',g)\mapsto\phi'(f')+\theta(g).
\]
我们有 $\delta'\alpha=0$。因此存在态射 $\delta\colon\Image\beta\to\cI$ 使 $\delta\beta=\delta'$。注意 $\Image\beta$ 是 $\cF$ 的含 $\cF'$ 的 $\cO_X$-子模，且 $\delta|_{\cF'}=\phi'$。另外，$\Image\beta$ 严格包含 $\cF'$，因为 $s$ 是 $\Image\beta$ 在 $W$ 上的截面，但不是 $\cF'$ 在 $W$ 上的截面。这与 $(\cF',\phi')$ 的极大性矛盾。所以 $\cF'=\cF$。这就证明了我们的断言。

设 $\cF$ 是 $\cO_X$-模，$I(\cF)$ 为所有从 $\cG$ 的 $\cO_X$-子模 $\cG_i$ 到 $\cF$ 的态射 $\phi_i\colon\cG_i\to\cF$ 的集合。考虑态射
\[
\phi\colon\bigoplus_{i\in I(\cF)}\cG_i\to\cF\oplus\cG^{|I(\cF)|},
\]
其中 $\cG^{|I(\cF)|}$ 是 $\cG$ 的若干份拷贝的直和，$I(\cF)$ 的每个元素对应其中一份拷贝 $\cG$；而 $\phi|_{\cG_i}\colon\cG_i\to\cF\oplus\cG^{|I(\cF)|}$ 是由 $-\phi_i\colon\cG_i\to\cF$ 与 $\cG_i$ 到 $\cG^{|I(\cF)|}$ 中第 $i$ 个直和项 $\cG$ 的典范包含所定义的态射。令 $\cI_1(\cF)$ 为 $\phi$ 的余核，$\phi(\cF)\colon\cF\to\cI_1(\cF)$ 为典范包含 $\cF\to\cF\oplus\cG^{|I(\cF)|}$ 与典范满射 $\cF\oplus\cG^{|I(\cF)|}\to\coker\phi$ 的复合。注意 $\phi(\cF)$ 是单态射。此外，每个态射 $\phi_i\colon\cG_i\to\cF$ 都可延拓为态射 $\cG\to\cI_1(\cF)$，后者由把 $\cG$ 嵌入 $\cF\oplus\cG^{|I(\cF)|}$ 中第 $i$ 份拷贝 $\cG$ 的态射 $\cG\to\cF\oplus\cG^{|I(\cF)|}$ 诱导。
\[\begin{tikzcd}[column sep=2.6em,row sep=1.5em]
\cG_i \arrow[r,"\phi_i"] \arrow[d] & \cF \arrow[d] \\
\cG \arrow[r] & \cI_1(\cF).
\end{tikzcd}\]

令 $\cS$ 为 $\cG$ 的所有 $\cO_X$-子模之集。选取良序集 $I$，使得对某个 $m\in I$，集合 $\{i\in I : i\le m\}$ 的基数严格大于 $|\cS|$。我们用超限归纳，对每个 $i\in I$ 定义 $\cI_i(\cF)$，并对每对 $j<i$ 定义单态射 $\cI_j(\cF)\to\cI_i(\cF)$，使得对每个三元组 $k<j<i$，下图交换：
\[\begin{tikzcd}[column sep=2.8em,row sep=1.6em]
\cI_k(\cF) \arrow[r] \arrow[dr] & \cI_j(\cF) \arrow[d] \\
& \cI_i(\cF).
\end{tikzcd}\]

令 $0$ 为 $I$ 的最小元。定义 $\cI_0(\cF)=\cF$。设 $i\in I$，并假设已对一切 $j<i$ 定义 $\cI_j(\cF)$，且对一切 $k<j<i$ 定义单态射 $\cI_k(\cF)\to\cI_j(\cF)$。若存在元素 $i-1\in I$，使没有 $i'\in I$ 具有性质 $i-1<i'<i$，则定义 $\cI_i(\cF)=\cI_1(\cI_{i-1}(\cF))$，并对 $j<i$ 把 $\cI_j(\cF)\to\cI_i(\cF)$ 定义为单态射 $\cI_j(\cF)\to\cI_{i-1}(\cF)$ 与单态射 $\cI_{i-1}(\cF)\to\cI_1(\cI_{i-1}(\cF))$ 的复合。若不存在这样的 $i-1$，则定义 $\cI_i(\cF)=\dlim_{j<i}\cI_j(\cF)$，并对 $j<i$ 把 $\cI_j(\cF)\to\dlim_{j<i}\cI_j(\cF)$ 定义为典范态射，它必为单态射。

设 $n$ 为 $I$ 中使
\[
|\{i\in I : i\le n\}|>|\cS|
\]
成立的最小元。我们有单态射 $\cF=\cI_0(\cF)\to\cI_n(\cF)$。我们来证明 $\cI_n(\cF)$ 是内射层。这将证明 $\cO_X$-模范畴有足够内射对象。只需证明：任一从 $\cG$ 的 $\cO_X$-子模 $\cG'$ 到 $\cI_n(\cF)$ 的态射 $\phi'$，其像含于某个 $n'<n$ 的 $\cI_{n'}(\cF)$。事实上，若这一点成立，则可把 $\phi'$ 诱导的态射 $\cG'\to\cI_{n'}(\cF)$ 延拓为态射 $\cG\to\cI_{n'+1}(\cF)$，其中 $n'+1$ 是 $I$ 中满足 $n'+1>n'$ 的最小元。把最后一个态射与单态射 $\cI_{n'+1}(\cF)\to\cI_n(\cF)$ 复合，得到态射 $\cG\to\cI_n(\cF)$，它延拓了态射 $\phi'\colon\cG'\to\cI_n(\cF)$。由开头的讨论，这说明 $\cI_n(\cF)$ 是内射层。

只需证明存在 $n'<n$，使对一切 $n'\le i<n$ 有
\[
\phi'^{-1}(\cI_i(\cF))=\phi'^{-1}(\cI_{n'}(\cF)).
\]
若不然，取 $J$ 为 $I$ 中满足 $j<n$ 且 $\phi'^{-1}(\cI_j(\cF))$ 严格含于 $\phi'^{-1}(\cI_{j+1}(\cF))$ 的元素 $j$ 的集合。我们有
\[
\{i\in I : i<n\}=\bigcup_{j\in J}\{i\in I : i\le j\}.
\]
每个 $\phi'^{-1}(\cI_j(\cF))$ 都是 $\cG$ 的 $\cO_X$-子模。所以 $|J|\le|\cS|$。由 $n$ 的最小性，对每个 $j\in J$ 有
\[
|\{i\in I : i\le j\}|\le|\cS|.
\]
于是
\[
\begin{aligned}
|\{i\in I : i<n\}|&=\Bigl|\bigcup_{j\in J}\{i\in I : i\le j\}\Bigr|\\
&\le\sum_{j\in J}|\cS|\\
&\le|J||\cS|\le|\cS|.
\end{aligned}
\]
这与 $|\{i\in I : i\le n\}|>|\cS|$ 矛盾。
\end{proof}
\begin{tnote}
2.1.6 的证明分三步：先用 Zorn 引理说明，只要对一个大对象 $\cG$ 的每个子模都有延拓性质，内射性即可得到；再对每个 $\cO_X$-模 $\cF$ 造出单态射 $\cF\to\cI_1(\cF)$，使 $\cG$ 的每个子模到 $\cF$ 的态射都能延拓到 $\cG$；最后用超限归纳与基数论证，把 $\cF$ 嵌入一个已具备该延拓性质的内射层 $\cI_n(\cF)$。论证只用到 $\cG$ 的子模集合 $\cS$ 的基数，故与 $X$ 的具体结构无关。[依据：原书 p.131--133]
\end{tnote}
